\chapter{MetaCIC TypedExamples module audit map}
\label{ch:metacic-typedexamples-audit-map}
\origin{ai}

\section{The setup}
\label{sec:metacic-typedexamples-audit-map-setup}

The BEDC.MetaCIC \texttt{TypedExamples} submodule is a worked-example gallery
for the internal CIC fragment. Its files collect Church encodings, dependent
identity terms, polymorphic identity terms, application forms, beta redexes,
projection forms, Pi forms, constants, and sort-indexed context examples.

This chapter reads the gallery as an audit surface. Unlike the subject
reduction, confluence, normalisation, substitution, and decidability audit maps,
the submodule is not primarily a metatheory module. It does not provide the
general proof that typing is preserved by every beta step or that every closed
typed term normalises. It supplies concrete typed terms whose elaboration and
typing witnesses exercise the substrate on representative programs.

The closure state is therefore gallery-level. Each cited target says that a
particular encoded term has the displayed type, closedness witness, beta
computation, or checker acceptance result. The collection is useful precisely
because the terms are ordinary objects of the represented syntax rather than
external examples.

\section{Inventory}
\label{sec:metacic-typedexamples-audit-map-inventory}

The module has four broad strata.

\begin{itemize}
\item \texttt{ChurchGallery.lean} and the \texttt{ChurchGallery/} directory
  provide the base Church encodings for booleans, natural numbers, pairs, and
  options. They expose the shared terms \texttt{churchBoolTy},
  \texttt{churchTrueTm}, \texttt{churchFalseTm}, \texttt{churchNatTy},
  \texttt{churchZeroTm}, \texttt{churchOneTm}, and \texttt{churchSuccTm}, plus
  typed constructors and eliminators for pair and option witnesses.
\item The sibling files \texttt{ChurchBoolLogic}, \texttt{ChurchNatRec},
  \texttt{ChurchNatArith}, \texttt{ChurchNatPred}, \texttt{ChurchEquality},
  \texttt{ChurchList}, and \texttt{ChurchSum} extend the gallery with boolean
  connectives, arithmetic operations, predecessor-style examples, Leibniz
  equality encodings, Church lists, and Church sums.
\item \texttt{Identity.lean} and \texttt{Polymorphic.lean} contain identity,
  dependent identity, polymorphic identity, constant, apply, flip, and compose
  forms. These examples check binder discipline and universe-like sort
  movement inside the represented calculus.
\item \texttt{AppForms.lean}, \texttt{Applications.lean},
  \texttt{ApplicationTail.lean}, \texttt{BasicApplications.lean},
  \texttt{BetaRedex.lean}, \texttt{ConstForms.lean}, \texttt{PiForms.lean},
  \texttt{Projections.lean}, and \texttt{Sorts.lean} cover application,
  beta-redex, constant, Pi, projection, and context lookup shapes.
\end{itemize}

This inventory is intentionally heterogeneous. Some files demonstrate direct
\texttt{HasType} derivations. Some use \texttt{inferTypeCtx\_sound} to turn a
checker result into a typing witness. Some record beta computation rows or
closedness rows that make later examples admissible. Together they form a
ground-truth set for the represented syntax.

The file-level audit reading is similarly explicit.

\begin{itemize}
\item The gallery root imports the base boolean, natural-number, pair, and
  option files and gives shared names to the recurring Church terms. Those
  shared names are the vocabulary consumed by the later arithmetic, equality,
  and list files.
\item \texttt{BoolNat.lean} is the smallest Church-encoding source. It proves
  the boolean and natural-number type formers, their constructors, and small
  identity and constant examples.
\item \texttt{Pair.lean} and \texttt{Option.lean} add product-like and
  maybe-like data. Their role is not merely to add more examples, but to force
  the checker through multi-argument constructors and eliminator-style terms.
\item \texttt{ChurchBoolLogic.lean} and the arithmetic files move from typing
  rows to computation rows. They record how boolean connectives and natural
  number operations reduce on selected inputs under the closedness hypotheses
  displayed in the theorem statements.
\item \texttt{ChurchEquality.lean} is the dependent proof-object branch. It
  checks Leibniz equality types, reflexivity proofs, closedness rows, and
  normal-equality checker rows for selected equality examples.
\item \texttt{ChurchList.lean} and \texttt{ChurchSum.lean} are the larger
  binder-stress files. They combine nested type formers, constructors,
  eliminators, closedness support, and computation rows for folds and case
  analysis.
\item The non-Church files remain important because they isolate syntax forms
  without the surrounding encoding burden. They are the local probes for sort
  contexts, Pi formation, projection behavior, constants, ordinary
  applications, application tails, and beta redexes.
\end{itemize}

Thus the inventory should not be read as a flat list of examples. It is a
layered gallery: atoms and binders first, then identity and application
patterns, then Church data, then eliminators and computation rows over that
data. The order matters for audit use because a later failure in a Church file
can often be compared against the smaller non-Church probes to locate whether
the issue is basic typing, binder shifting, eliminator typing, or beta
computation.

\section{Representative theorems}
\label{sec:metacic-typedexamples-audit-map-representative-theorems}

The base Church boolean and Church natural-number rows are the smallest public
examples. The boolean type and constructors are typed in the empty context:
\leanvariant{BEDC.MetaCIC.church\_bool\_type}
\leanvariant{BEDC.MetaCIC.church\_true}
\leanvariant{BEDC.MetaCIC.church\_false}.
The natural-number encoding and its first constructors are likewise typed:
\leanvariant{BEDC.MetaCIC.church\_nat\_type}
\leanvariant{BEDC.MetaCIC.church\_zero}
\leanvariant{BEDC.MetaCIC.church\_one}
\leanvariant{BEDC.MetaCIC.church\_succ}.
These targets establish that the gallery's basic numerals and booleans are not
paper-only conventions; they are terms accepted by the encoded typing relation.

The list gallery checks a type former, constructors, and a fold operator:
\leanvariant{BEDC.MetaCIC.church\_list\_type\_typed}
\leanvariant{BEDC.MetaCIC.church\_nil\_typed}
\leanvariant{BEDC.MetaCIC.church\_cons\_typed}
\leanvariant{BEDC.MetaCIC.church\_fold\_typed}.
This group exercises nested Pi structure, shifted element types, and eliminator
typing. It is a compact stress test for binder indices because the encoded list
type moves an element type through several binders.

The sum gallery records the either type, both injections, and case analysis:
\leanvariant{BEDC.MetaCIC.church\_either\_type\_typed}
\leanvariant{BEDC.MetaCIC.church\_inl\_typed}
\leanvariant{BEDC.MetaCIC.church\_inr\_typed}
\leanvariant{BEDC.MetaCIC.church\_case\_either\_typed}.
The accompanying beta rows, such as
\leanvariant{BEDC.MetaCIC.church\_case\_inl}
\leanvariant{BEDC.MetaCIC.church\_case\_inr},
show that the eliminator computes on both injection branches under the
displayed closedness inputs.

The equality gallery uses a Leibniz-style encoding:
\leanvariant{BEDC.MetaCIC.leibniz\_eq\_type}
\leanvariant{BEDC.MetaCIC.leibniz\_refl\_type}
\leanvariant{BEDC.MetaCIC.refl\_zero\_proof}
\leanvariant{BEDC.MetaCIC.refl\_true\_proof}.
These rows show a dependent proof object, not merely first-order data. They are
useful witnesses that the same typed-example layer includes proposition-shaped
programs.

The dependent and polymorphic identity examples give a separate binder-focused
surface:
\leanvariant{BEDC.MetaCIC.lam\_dependent\_identity}
\leanvariant{BEDC.MetaCIC.polymorphic\_identity}
\leanvariant{BEDC.MetaCIC.poly\_identity\_type}
\leanvariant{BEDC.MetaCIC.poly\_identity\_at\_sort\_type}.
They test the common identity pattern in both direct dependent and explicitly
polymorphic forms.

The application and beta-redex files make the gallery operational rather than
only declarative:
\leanvariant{BEDC.MetaCIC.app\_id\_sort\_in\_empty\_ctx}
\leanvariant{BEDC.MetaCIC.dep\_id\_applied\_to\_arg\_in\_sort\_ctx}
\leanvariant{BEDC.MetaCIC.id\_sort\_redex}
\leanvariant{BEDC.MetaCIC.id\_sort\_redex\_typed}.
These examples connect the typing gallery to the beta relation used elsewhere
in MetaCIC.

\section{Role within MetaCIC}
\label{sec:metacic-typedexamples-audit-map-role}

The submodule functions as the typing relation's ground-truth gallery. A change
to the typing implementation, closedness infrastructure, substitution helpers,
or beta evaluator should preserve these examples unless the represented
calculus itself has changed. In that sense the gallery is a regression surface
for concrete syntax, binder shifting, context lookup, application typing, and
closed Church-encoded programs.

The gallery also keeps the manuscript honest about examples. When the paper
mentions Church booleans, naturals, lists, sums, equality, or identity programs
inside MetaCIC, the cited rows above are the Lean-side witnesses that the terms
exist with the stated typing or computation behavior.

\section{Connection to other audit maps}
\label{sec:metacic-typedexamples-audit-map-connection}

The closest neighbour is the decidability map in
\autoref{ch:metacic-decidable-audit-map}. Its \texttt{GalleryTests} layer checks
that the type inference surface accepts the main closed example families. The
present chapter records the worked examples themselves. The two layers should
be read together: \texttt{TypedExamples} supplies the concrete programs, while
the decidability gallery checks the option-valued inference interface on those
program families.

The subject-reduction map in
\autoref{ch:metacic-subject-reduction-obstruction} has a different role.
TypedExample terms that take beta steps should preserve typing along the
expected route, but the audit map cannot promote that expectation into an
unconditional subject-reduction theorem. The SR chapter names the preservation
obligations that would be needed for that stronger claim.

The normalisation and confluence maps are also adjacent but not identical. The
Church examples include concrete beta computations and some normalisation-style
rows. They are witnesses that the reduction machinery runs on selected terms,
not a proof of arbitrary closed-term normalisation or arbitrary closed-term
Church--Rosser.

\section{Boundary}
\label{sec:metacic-typedexamples-audit-map-boundary}

This chapter is an audit map for a worked-example collection. It may be cited
for the displayed typed Church encodings, dependent identity examples,
polymorphic identity examples, application examples, beta-redex rows, and
selected computation rows.

It does not prove subject reduction, strong normalisation, confluence, or global
conversion decidability. It does not claim that every Church-style encoding has
all expected algebraic laws inside MetaCIC. It records the examples whose Lean
targets are present and keeps their role at gallery level.

\concretizedIn{ch:concrete-instances-metacictypedexamplesgallery-namecert}{TypedExamples gallery rows are represented as a finite NameCert regression surface.}
\concretizedIn{ch:concrete-instances-typedexamplegallery-namecert}{Typed example gallery rows are packaged as a finite NameCert carrier with checker, closedness, beta, and regression-surface obligations.}
\concretizedIn{ch:concrete-instances-metacic-typedexamples-seal-namecert}{TypedExamples gallery rows are realized as a finite worked-example seal.}
\concretizedIn{ch:concrete-instances-typedexampleauditbudget-namecert}{Typed example audit budgets realize the TypedExamples audit map as finite regression-spend NameCert rows.}
\concretizedIn{ch:concrete-instances-metacictypedexamplesapplicationroute-namecert}{Application-form examples are realized as finite typed-application route packets.}
\concretizedIn{ch:concrete-instances-typedexamplecoveragematrix-namecert}{Typed example coverage matrices realize gallery, checker, beta, and regression-spend coverage as finite NameCert rows.}
\concretizedIn{ch:concrete-instances-church-eliminator-regression-matrix-namecert}{Church eliminator regression matrices realize bool, nat, list, and sum eliminator gallery rows as one finite checker and beta-regression NameCert packet.}
\concretizedIn{ch:concrete-instances-typed-example-reduction-trace-namecert}{Typed-example reduction traces are realized as finite MetaCIC NameCert rows linking typed term, reduction step, checker result, closedness, transport, replay, and provenance.}
